\begin{titlepage}\centering
\vspace*{1cm}{\Large Carl Friedrich Gauss\par}\bigskip
{\huge Arithmetische Abhandlungen\par}\bigskip
{\Large Deutsche Übersetzung\par}\medskip
{\large Band II: arithmetische Abhandlungen, Besprechungen und Nachlassstudien\par}\vfill
\begin{minipage}{0.90\linewidth}\small
Diese Auswahl enthält die historischen Titelblätter und die ersten fünf vollständigen Abhandlungen des zweiten Bandes: Einen neuen Beweis eines arithmetischen Satzes (Artikel 1–7); die Summation einiger besonderer Reihen (Artikel 1–36); Neue Beweise und Erweiterungen des Fundamentaltheorems in der Theorie der quadratischen Reste, mit der Einleitung, dem fünften Beweis (Artikel 1–3), dem sechsten Beweis (Artikel 1–6) und dem neuen Algorithmus (Artikel 1–6); sowie beide Abhandlungen zur Theorie der biquadratischen Reste (Artikel 1–23 und 24–76). Es folgen fünf Selbstbesprechungen von 1808–1831 und sieben Besprechungen fremder Schriften von 1809–1831. Letztere behandeln irrationale Wurzelgrößen, die Faktorentafeln von Chernac und Burckhardt, Erchingers Siebzehneckkonstruktion und Seebers positive ternäre quadratische Formen.

Die Nachlassstudien enthalten die Lösung einer reinen Kongruenz, Lehre von den Resten, sechstes Kapitel, erster Teil (Artikel 237–252); die Allgemeinen Untersuchungen über Kongruenzen, achtes Kapitel (Artikel 330–375 wie gedruckt, einschließlich des unvollendeten Schlusses von Artikel 375), mit R. Dedekinds historischen Bemerkungen; und die Weitere Entwicklung der Untersuchungen über reine Gleichungen, vollständig wie gedruckt (Artikel 1–19), mit Dedekinds Bemerkungen. Es folgt das Fragment Beweis einiger Sätze über die Perioden der Klassen binärer quadratischer Formen, mit beiden Sätzen, dem unvollendeten Schluss des zweiten Beweises und Sterns historischer Bemerkung.

Die mit [I.] bezeichnete Fassung der ersten Abhandlung Über den Zusammenhang zwischen der Anzahl der Klassen, in die binäre quadratische Formen zerfallen, und ihrer Determinante folgt vollständig wie gedruckt, mit dem Titel und den Artikeln 1–8.

Die ursprünglich deutschen Titelangaben, Besprechungen, der Abteilungstitel des Nachlasses und die historischen Bemerkungen sind in modernem Deutsch wiedergegeben. Ihre historischen Originale werden gesondert mitgeliefert. Spätere Texte des zweiten Bandes liegen außerhalb dieser Auswahl. Der vollständige erste Band ist im gesammelten Leseband enthalten.

Historische Vorlage: Carl Friedrich Gauss, Werke, zweiter Band, zweiter Abdruck, Göttingen, 1876.

Deutsche Übersetzung erstellt mit OpenAI GPT-6.1 Sol.
\end{minipage}\end{titlepage}
